% 第 66 章骨架：导引框已定，正文待写。写作计划见 docs/PLAN.md“第 66 章”，写法见 docs/STYLE.md。
\chapter{零利率下限、量化宽松与前瞻指引}\label{ch:zlb}

\begin{chapterguide}
\guidequestion{利率降到零以后，央行还能做什么？}
\begin{guidegoals}
\item 用新凯恩斯模型分析零利率下限与流动性陷阱（Krugman 1998；Eggertsson \& Woodford 2003）；
\item 解释量化宽松的作用渠道与实证估计；
\item 说明前瞻指引与“前瞻指引之谜”；
\item 比较美联储、英格兰银行、欧洲央行与日本银行的做法。
\end{guidegoals}
\guideprereq{\chref{ch:consensus}\enspace\nameref{ch:consensus}；\chref{ch:japanbubble}\enspace\nameref{ch:japanbubble}}
\end{chapterguide}

\chapterstub
